\documentclass[11pt]{article}
\usepackage[a4paper,margin=25mm]{geometry}
\usepackage{amsmath,amssymb,amsthm,booktabs,hyperref}
\newtheorem{theorem}{Theorem}
\newtheorem{lemma}{Lemma}
\newcommand{\Conf}{\mathsf{Conf}}
\newcommand{\Halt}{\mathsf{Halt}}
\title{A short tape detour: from 243 to 242 states}
\author{Research contribution prepared for the ProveIt project}
\date{18 October 2026}
\begin{document}
\maketitle

\paragraph{Scope.}
This note gives an ordinary mathematical proof of the exact binary-machine
transformation for arbitrary head positions and tapes at the shared start
state 95. A small structural invariant proves that old 27 always reads 0;
no source-reachability hypothesis or bounded execution experiment is
required. The accompanying checker verifies the exact table patch, every
incoming row needed by that invariant, all whole-tape block identities,
and the complete finite classification of transition rows. It does not
claim a Lean kernel theorem or independently prove source calibration.

\section{Exact transformation}

All state numbers in this note are the inherited state labels, not the
canonical consecutive labels. Let $M$ be Combined243 with start 95 and
halt 254. Its inherited table has SHA256
\begin{quote}\ttfamily\small
4207fc3d1bb16f3c64c3cff0980eaa1486160c3f467aee7685740a39e3df7300.
\end{quote}
Let $N$ be obtained by deleting state 27, redirecting every incoming edge of 27
to 15, and replacing $(15,0)$ by $0L32$. The complete list of changed
retained rows is
\begin{center}
\begin{tabular}{rrrr}
\toprule
State&Read&Old action&New action\\
\midrule
15&0&$0R21$&$0L32$\\
18&0&$0L27$&$0L15$\\
47&0&$0L27$&$0L15$\\
65&0&$0L27$&$0L15$\\
\bottomrule
\end{tabular}
\end{center}
The two rows of 27 disappear. The other 480 retained rows remain literally
unchanged. Consequently $N$ has 242 working states, 484 total binary rows,
and one separate halt state. All its transitions move left or right. No
transition targets 27. The inherited candidate SHA256 is
\begin{quote}\ttfamily\small
c8ef2f44d25fee6813164cc59ec05061799397da8cf7f0ea7baab63d32d76e02.
\end{quote}
The canonical table numbers the start as 0, the other working states as
1 through 241, and halt as 242. Its SHA256 is
\begin{quote}\ttfamily\small
ef8c07f372c76a73c021f51801f16f28f1e4ca5ead2d4ce6780c2a245a0db245.
\end{quote}

\section{A one-cell invariant}

A configuration is a triple $C=(q,h,\tau)$ with $q$ a control state,
$h\in\mathbb Z$, and $\tau:\mathbb Z\to\{0,1\}$ the entire tape.
Put $S=\{15,18,47,65\}$ and define
\[
 I(C)\quad\Longleftrightarrow\quad
 \bigl(q\in S\Longrightarrow\tau(h-1)=0\bigr)
 \;\land\;
 \bigl(q=27\Longrightarrow\tau(h)=0\bigr).
\]

\begin{lemma}
Every configuration in an $M$ execution starting in 95 satisfies $I$,
regardless of its initial head position and tape.
\end{lemma}
\begin{proof}
The initial control is outside $S\cup\{27\}$. The complete incoming-row
census for the states in $S$ is
\begin{center}
\begin{tabular}{rll}
\toprule
Target&Source rows&Action\\
\midrule
15&$(11,0)$&$0R15$\\
18&$(12,0),(148,0)$&$0R18$\\
47&$(33,0),(144,0),(167,0)$&$0R47$\\
65&$(51,0),(81,0)$&$0R65$\\
\bottomrule
\end{tabular}
\end{center}
Every one of these rows writes 0 and moves right, proving the left-zero
property at the new head. The only rows targeting 27 are
$(18,0),(47,0),(65,0)$, each with action $0L27$. Their source controls
belong to $S$, so the inductive hypothesis says that the cell one position
left of the predecessor head is 0. The write takes place at the predecessor
head and does not change that left cell. Moving left therefore enters 27
on 0. At every other target both implications are vacuous. Induction on
the execution length proves the entire invariant, including the
unreachability of row $(27,1)$.
\end{proof}

\section{Whole-tape comparison blocks}

Write $\xrightarrow[k]{M}$ for exactly $k$ transitions. The changed row
at 15 is justified by the following identity, not by identifying 15 and 27
as interchangeable states on arbitrary tapes.

\begin{lemma}[The detour]
If $\tau(h)=\tau(h-1)=0$, then
\[
 (15,h,\tau)\xrightarrow[1]{M}(21,h+1,\tau),\qquad
 (15,h,\tau)\xrightarrow[3]{N}(21,h+1,\tau).
\]
If $\tau(h)=1$, both machines instead take one identical step to
$(22,h+1,\tau)$.
\end{lemma}
\begin{proof}
The three new rows are exactly
\[
 (15,0)\mapsto(0,L,32),\quad
 (32,0)\mapsto(0,R,45),\quad
 (45,0)\mapsto(0,R,21).
\]
They visit head positions $h,h-1,h$. The inspected cells are both 0;
every write preserves the inspected cell. The final state, head, and
entire tape therefore agree with the old one-step result. On input 1,
$(15,1)\mapsto(1,R,22)$ was not changed.
\end{proof}

\begin{lemma}[Redirected entries]
For $p\in\{18,47,65\}$, if $\tau(h)=\tau(h-1)=0$, then
\[
 (p,h,\tau)\xrightarrow[2]{M}(32,h-2,\tau),\qquad
 (p,h,\tau)\xrightarrow[2]{N}(32,h-2,\tau).
\]
\end{lemma}
\begin{proof}
The old first step reaches 27 at $h-1$, while the new first step reaches 15
there. Their respective zero rows are both $0L32$. Both cells inspected
were assumed 0, so every write preserves the complete tape.
\end{proof}

In these lemmas every unspecified cell of $\tau$ is arbitrary. The proofs
neither assume that cells outside a finite window are zero nor discard
changes outside an observed region. The finite checker represents each
unread exterior cell by a distinct formal atom, and compares normalized
whole-tape overlays.

\section{Global relative halting theorem}

Define $f(27)=15$ and $f(q)=q$ for every other inherited control,
including halt 254. Exact finite inspection gives
\[
 \delta_N(f(q),b)=(w,d,f(t))
 \quad\text{when}\quad\delta_M(q,b)=(w,d,t)
\]
for all 484 old rows other than $(15,0)$ and $(27,1)$. This includes the
three redirected incoming rows and the old row $(27,0)$:
old 27 on 0 and new15 on 0 both take the identical step $0L32$.
The row $(27,1)$ is inadmissible by the structural invariant.

\begin{theorem}
Let $C_0=(95,h_0,\tau_0)$, with arbitrary head $h_0$ and tape $\tau_0$.
Then
\[
 \Halt(M,C_0)\quad\Longleftrightarrow\quad\Halt(N,C_0).
\]
\end{theorem}
\begin{proof}
Relate the old configuration $(q,h,\tau)$ to the new configuration
$(f(q),h,\tau)$ whenever the old configuration satisfies $I$.
The old structural invariant applies at every step. Since $f(95)=95$,
the two initial configurations are related.

For the old row $(15,0)$, the invariant supplies a zero immediately
left. The detour lemma simulates the old step by three new steps to the
same configuration, whose state 21 is fixed by $f$. Every other
admissible old row is simulated by one new step by the mapped-row
identity just displayed. The apparent overlap at new15 is harmless:
an old 27 configuration uses the one-step identity to 32, while an old 15
configuration on 0 uses the full three-step detour to 21.

Every old step therefore corresponds to one or three new steps. If the
old execution has taken $n$ steps, of which $k$ used $(15,0)$, the related
new configuration occurs at exactly $n+2k$ steps. Both head and entire
tape agree, and the controls correspond by $f$. In particular, the
new elapsed time lies between $n$ and $3n$ and tends to infinity with
$n$. The two proper intermediate controls in the only longer block are 32
and 45, both working controls. Finally $f$ maps the old halt to the new
halt and maps every working control to a working control.

Consequently an old finite halt is simulated by a new finite halt, and
an infinite old execution supplies an infinite new execution. A new
halt cannot occur inside the longer block, so a new finite halt must
also correspond to an old finite halt. This establishes the claimed
equivalence without a consistency assumption.
\end{proof}

\section{Consistency calibration}

The structural invariant closes the potential $(27,1)$ issue directly
at the binary-machine level. In particular, the comparison theorem
applies to blank input without any additional hypothesis about the
source program, reachable register values, or gap size. Composing it
with the existing ordinary Combined243 calibration gives
\[
 \mathrm{ZFC}\vdash
 \bigl(\neg\Halt(N,\mathrm{blank})\leftrightarrow
       \operatorname{Con}(\mathrm{ZFC})\bigr).
\]
The local arguments use finite transition calculations and ordinary
induction, so they can be formalized in ZFC. The bootstrap and the fresh
full-frame replay of the 242-state table are useful independent checks,
but are not premises of the relative machine theorem. This statement
distinguishes an ordinary proof supported by executable certificates
from completion of a new Lean consistency endpoint. Existing gaps in
the predecessor's full Lean consistency bridge remain separate from
the ordinary binary reduction established here.

\paragraph{Attribution.}
The predecessor calculus, program, compiler, and reductions retain their
existing attributions to Mira-Cult, Jeff Ketchersid, Andrew James Wade,
and the ProveIt and incoming research contributions. This contribution
is the additional 27-to 15 detour and its exact local comparison proof.

\end{document}
